The implementation uses CUDA's SIMT execution: threads are organised into
blocks and 32-thread warps, with registers, block-shared memory and device
global memory serving different reuse scopes \cite{cuda_programming_guide}.
Coalescing groups neighbouring lane accesses into suitable memory
transactions; it does not require fewer logical code bits. Candidate-parallel
primary scanning assigns lanes to neighbouring vectors, while refinement
assigns several lanes to one survivor's coordinates. These mappings exploit
different available parallelism, so one mapping need not suit both stages.

Occupancy is constrained by registers, shared state and block geometry.
More resident threads can hide latency, but more blocks per query can also
duplicate tables and local candidate buffers. Shared memory is valuable
when query data are reused enough to repay loading and synchronisation.
These observations motivate Chapter 5's table grouping, packed access and
batch-aware scheduling; the throughput ablations assess their combined
effects. A claim of bandwidth saturation or a specific cache bottleneck
would additionally require counters and a consistent physical-byte model.

Selection is itself substantial computation. Faiss's GPU work combines
efficient distance processing with warp-oriented selection
\cite{johnson_gpu}. Bitonic networks, radix methods and merge-based
approaches provide established building blocks with different $k$, input
length and state costs. AIR top-$k$ and GridSelect further illustrate why
global selection design depends on workload and whether scores arrive
on the fly \cite{gpu_topk}. This thesis adapts such ideas to JHQ's probe
selection and survivor retention; it does not claim to invent these primitives.

CUDA Graph replay reduces repeated launch setup for a captured workload.
Capture, workspace allocation and changes in graph shape belong to the
index/search lifecycle, rather than the mathematical scoring rule.
Likewise, publishing a shared threshold requires an explicit ordering and
memory-access contract. Correct scalar scores do not alone prove concurrent
admission safe. Chapter 5 separates these correctness conditions from
performance measurements, and Chapter 3 keeps device profiling distinct
from the host-visible search interval.
